\subsection{函子 $J\rightarrow
\Ind(C)$ 的带指标表示}\etiquette[8.8]{I.sec8.8}\pageoriginale

\subsubsection{}\etiquette[8.8.1]{I.sec8.8.1}
设
$$
\varphi:J\longrightarrow \Ind(C)\text{ \ie } \varphi\in \Ob
\Hhom(J,\Ind(C))
$$
为一个函子, 其中 $C$ 是 $𝒰$-范畴. 回忆\hyperref[I.sec8.5.4]{(8.5.4)},
按定义, $\varphi$ 的一个\emph{带指标表示}是 $\Hhom(J,C)$ 的一个 ind-对象
$$
\psi:I\longrightarrow \Hhom(J,C){\quoi},
$$
其中 I 是本质小的滤过范畴, 并给定其在 $\Hhom(J,\Ind(C))$ 中的归纳极限到 $\varphi$ 的一个同构.
因此, $\varphi$ 有带指标表示, 当且仅当它同构于 $\Hhom(J,\Ind(C))$
中由全子范畴 $\Hhom(J,C)$ 的对象组成的一个本质小的滤过归纳极限.
若每个函子 $\varphi:J\rightarrow \Ind(C)$ 都有带指标表示,
则称范畴 $J$ \emph{对 $C$ 可容许}.

由于在 $\Ind(C)$ 中小的滤过归纳极限可表
(\hyperref[I.prop8.5.1]{8.5.1}), 在
$\Hhom(J,\Ind(C))$ 中也如此, 而且这些极限是“逐个变元”计算的.
因此, 包含函子
\hypertarget{sga4.exp.I.equation.8.8.1.1}{}$$
\displaylines{\rlap{$(8.8.1.1)$}\hfill \Hhom(J,C)\hookrightarrow
\Hhom(J,\Ind(C))\hfill}
$$
典范地延拓为一个函子
\hypertarget{sga4.exp.I.equation.8.8.1.2}{}$$
\displaylines{\rlap{$(8.8.1.2)$}\hfill
\Ind(\Hhom(J,C))\longrightarrow \Hhom(J,\Ind(C))\hfill}
$$
\hyperref[I.sec8.7.2]{(8.7.2)}. 这个函子的本质像的对象恰好是有带指标表示的那些 $\varphi$;
因此, 说 $J$ 对 $C$ 可容许, 也就是说函子\hyperlink{sga4.exp.I.equation.8.8.1.2}{(8.8.1.2)}
本质满射. 关于这一点, 我们指出:

\begin{enonce*}{命题 8.8.2}\etiquette[8.8.2]{I.prop8.8.2}
若\pageoriginale 范畴 $J$ 等价于一个有限范畴,
则典范函子 $\hyperlink{sga4.exp.I.equation.8.8.1.2}{(8.8.1.2)}$ 全忠实;
因而 $J$ 对 $C$ 可容许, 当且仅当这个函子是范畴等价.
\end{enonce*}

由\hyperref[I.prop8.7.5]{8.7.5} a), 只须证明, 对 $\Hhom(J,C)$ 中每个小的滤过归纳系
$(\varphi_i)_{i\in I}$ 及 $\Hhom(C,J)$ 的每个对象 $\varphi$, 典范映射
\hypertarget{sga4.exp.I.equation.8.8.2.1}{}$$
\displaylines{\rlap{$(8.8.2.1)$}\hfill \varinjlim_i \Hom(\varphi,
\varphi_i)\longrightarrow \Hom(\varphi, \varinjlim_i
\varphi_i)\hfill}
$$
是双射, 其中 $\varinjlim_i \varphi_i$ 表示在 $\Hom(J,\Ind(C))$ 中取的归纳极限.
若 $\varphi$, $\psi$ 是两个函子 $J\rightarrow C$, 则有一个集合的正合图,
且它关于 $\varphi$ 和 $\psi$ 具有函子性:
$$
\Hom(\varphi,\psi)\longrightarrow \prod_{X\in \Ob J}
\Hom(\varphi(X),\psi(X))\rightrightarrows \prod_{\substack{f\in
F\ell J\\ f:X\rightarrow Y}} \Hom(\varphi(X), \psi (Y)){\quoi}.
$$
由于滤过归纳极限与双箭头的核及有限乘积可交换,
当 $J$ 有限时便得到\hyperlink{sga4.exp.I.equation.8.8.2.1}{(8.8.2.1)} 的双射性.
$J$ 等价于一个有限范畴的情形立即归结为前一种情形.

\begin{enonce*}{命题 8.8.3}\etiquette[8.8.3]{I.prop8.8.3}
设 $\varphi:J\rightarrow \Ind(C)$ 为函子, $J$ 等价于一个小范畴.
下列两个条件足以保证 $\varphi$ 有带指标表示;
当 $J$ 等价于一个有限范畴时, 这两个条件也是必要的:
\begin{enumerate}
\renewcommand{\theenumi}{\alph{enumi}}
\renewcommand{\labelenumi}{\theenumi)}
\item 范畴 $\Hhom(J,C)_{/\varphi}$ (由 $\Hhom(J,\Ind(C))$ 中终点为 $\varphi$、
起点在 $\Hom(J,C)$ 中的箭头组成) 是滤过的.

\item 对\pageoriginale $J$ 的每个对象 $j$, 函子
\hypertarget{sga4.exp.I.equation.8.8.3.1}{}$$
\displaylines{\rlap{$(8.8.3.1)$}\hfill \psi\longmapsto
\psi(j):\Hhom(J,C)_{/\varphi}\longrightarrow C_{/\varphi(j)}\hfill}
$$
\end{enumerate}
是共尾的, \ie 满足 $\hyperref[I.prop8.1.3]{8.1.3}$ 的条件 {\rm %%
F 1), F 2)}.
\end{enonce*}

\noindent \emph{证明}. 设 a) 和 b) 成立;
来证明 $\varphi$ 有带指标表示. 显然可设 $J$ 是小的. 令
\hypertarget{sga4.exp.I.equation.8.8.3.2}{}$$
\displaylines{\rlap{$(8.8.3.2)$}\hfill
I=\Hhom(J,C)_{/\varphi}{\quoi},\hfill}
$$
由假设它是滤过范畴. 先设 $I$ 本质小, 于是 $I$ 到 $\Hhom(J,C)$ 的“包含”
定义 $\Hhom(J,C)$ 的一个 ind-对象 $i\mapsto
\psi_i$.
此外还有一个典范同态
\hypertarget{sga4.exp.I.equation.8.8.3.3}{}$$
\displaylines{\rlap{$(8.8.3.3)$}\hfill \varinjlim_i
\psi_i\longrightarrow \varphi{\quoi},\hfill}
$$
它逐个变元由函子\hyperlink{sga4.exp.I.equation.8.8.3.1}{(8.8.3.1)} 诱导的同态
$$
\varinjlim_I \psi_i (j)\longrightarrow
\varphi(j)=\varinjlim_{C_{/\varphi(j)}}X
$$
给出. 由于%%
后一个函子共尾, 得知\hyperlink{sga4.exp.I.equation.8.8.3.3}{(8.8.3.3)} 是同构,
这就证明了断言. 当不假定 $I$ 本质小时, 只须构造一个本质小的全子范畴 $I'$,
使得将 $\varinjlim_I$ 换成 $\varinjlim_{I'}$ 后,
\hyperlink{sga4.exp.I.equation.8.8.3.3}{(8.8.3.3)} 仍是同构.
为此, 只须由函子\hyperlink{sga4.exp.I.equation.8.8.3.1}{(8.8.3.1)}
诱导\pageoriginale 的复合函子
$$
I'\longrightarrow I\longrightarrow C_{/\varphi(j)}
$$
全部共尾. 于是由下列引理即可得出结论;
该引理的证明由判别准则\hyperref[I.prop8.1.3]{(8.1.3 b)}) 立即得出, 留给读者:

\begin{enonce*}{引理 8.8.3.4}\etiquette[8.8.3.4]{I.lem8.8.3.4}
设 $I$ 是滤过 $𝒰$-范畴, 且
$$
f_j:I\longrightarrow I_j\quad (j\in J)
$$
是从 $I$ 到本质小的滤过范畴的一族小的共尾函子.
则 $I$ 存在全子范畴 $I'$, 它滤过且本质小,
并且由各 $f_j$ 诱导的函子共尾.
\end{enonce*}

最后, 来证明当 $J$ 等价于有限范畴时\hyperref[I.prop8.8.3]{8.8.3} 中的必要性.
用本质小的滤过指标范畴 $I$ 给出的 $\varphi$ 的一个带指标表示定义一个函子 $\psi$
\hypertarget{sga4.exp.I.equation.8.8.3.5}{}\[
\xymatrix@C=1.25cm@R=2cm{
 I \ar[rr]^(.4){\psi} \ar@{-->}[rrd]^{\psi_j}&&  \Hhom(J,C)_{/\varphi} \ar@{-->}[d] &\\
                        &     &   C_{/\varphi(j)}\quad ,\\
}\tag{8.8.3.5}
\]
只须证明 $\psi$ 以及它所诱导的每个函子 $\psi_j$ 都共尾:
由\hyperref[I.prop8.1.3]{8.1.3} b), 这将推出 $\Hhom(J,C)_{/\varphi}$ 滤过,
而由定义\hyperref[I.def8.1.1]{8.1.1},
\hyperlink{sga4.exp.I.equation.8.8.3.5}{(8.8.3.5)} 中的竖直箭头也共尾.
由\hyperref[I.prop8.8.2]{8.8.2}, 可以把 $\varphi$ 等同于 $\Ind(E)$ 的一个对象,
其中 $E=\Hhom(J,C)$. 函子
$$
\psi:I\longrightarrow E_{/\varphi}{\quoi},
$$
由 $E$ 的以 $I$ 为指标的 ind-对象 $\varphi$ 产生\pageoriginale;
这个函子共尾是一个一般事实, 由\hyperref[I.prop8.1.3]{8.1.3} 的判别准则 (F 1) 和 (F 2) 立即得到.
同一论证 (将 $E$, $\varphi$ 换成 $C$, $\varphi(j)$)
表明 $\psi_j$ 共尾, 这就完成了证明.

\begin{enonce*}[remark]{注 8.8.4}\etiquette[8.8.4]{I.rem8.8.4}
\begin{enumerate}
\renewcommand{\theenumi}{\alph{enumi}}
\renewcommand{\labelenumi}{\theenumi)}
\item 当 $J$ 等价于一个有限范畴时, 若 $\varphi$ 有带指标表示,
我们已经看到\hyperlink{sga4.exp.I.equation.8.8.3.5}{(8.8.3.5)} 是共尾函子,
这蕴含范畴 $\Hhom(J,C)_{/\varphi}$ 本身本质小.

\item 设 $C$ 中有限归纳极限可表.
则 $E=\Hhom(J,C)$ 中也如此, 且这些极限逐个变元计算;
它们也是 $\Hhom(J,\widehat{C})$ 中的归纳极限,
\textit{从而更是} $\Hhom(J,\Ind(C))$ 中的归纳极限.
立即得知, 此时\hyperref[I.prop8.8.3]{8.8.3} 的条件 a) 自动满足,
只剩下条件 b) 需要考虑. 我不知道在进一步假定 $J$ 是小时, 该条件是否自动满足.
\end{enumerate}
\end{enonce*}
现在给出本节的主要结果:

\begin{enonce*}{命题 8.8.5}\etiquette[8.8.5]{I.prop8.8.5}
设范畴 $J$ 等价于一个有限范畴, 并进一步假定 $J$ 是刚性的,
\ie 对每个 $j\in \Ob J$, $j$ 的每个自同态都是恒等态射.
则 $J$ 对 $C$ 可容许 (其中 $C$ 是任意 $𝒰$-范畴),
\ie 每个函子 $J\rightarrow \Ind(C)$ 都有带指标表示.
\end{enonce*}

必要时将 $J$ 换成等价范畴\pageoriginale, 可设 $J$ 是\emph{骨架式}的,
\ie $J$ 中两个同构的对象必相等. 此时 $J$ 实际上有限.
对 $\text{card}\Ob J$ 作归纳; 该数 $\leq 0$ 的情形平凡, 因而假定它 $\geq 1$.
取 $J$ 的一个\emph{极大}对象 $j_\circ$, \ie 对每个箭头 $j_\circ\rightarrow j$,
存在箭头 $j\rightarrow j_\circ$; 因为 $J$ 是刚性的, 这蕴含 $j_\circ\rightarrow j$ 是同构,
又因 $J$ 为骨架式, 得到 $j_\circ=j$.
于是令 $J'$ 为从 $J$ 去掉对象 $j_\circ$ 得到的全子范畴,
令 $J''$ 为仅由对象 $j_\circ$ 组成的全子范畴.
该命题便由下列引理得到:

\begin{enonce*}{引理 8.8.5.1}\etiquette[8.8.5.1]{I.lem8.8.5.1}
设 $J$ 是有限范畴, $J'$ 和 $J''$ 是两个全子范畴,
使得 $\Ob J=\Ob J'\cup \Ob J''$, 且对每个 $j'\in \Ob J'$ 和 $J''\in \Ob J''$,
都有 $\Hom(j'',j')=∅$. 若 $J'$ 和 $J''$ 对 $C$ 可容许, 则 $J$ 亦然.
\end{enonce*}

只须验证\hyperref[I.prop8.8.3]{8.8.3} 的判别条件 a) 和 b),
这相当于作六项初等验证, 即 $\Hhom(J,C)_{/\varphi}$ 的后两个滤过性条件,
以及对函子\hyperlink{sga4.exp.I.equation.8.8.3.1}{(8.8.3.1)},
依次在 $j\in \Ob J'$ 与 $j\in \Ob J''$ 两种情形下验证条件 F 1) 和 F 2);
后面的条件还蕴含 $\Hhom(J,C)_{/\varphi}$ 非空.
这一验证颇为繁琐, 但没有困难, 留给读者.
(编者曾试图找到一个优美的证明, 以绕开这些令人不快的验证, 但未能如愿.)

\begin{enonce*}[remark]{习题 8.8.6}\etiquette[8.8.6]{I.exer8.8.6}
\begin{enumerate}
\renewcommand{\theenumi}{\alph{enumi}}
\renewcommand{\labelenumi}{\theenumi)}
\item 设\pageoriginale $J$, $J'$ 是两个对 $C$ 可容许的范畴, 其中一个有限.
证明 $J\times J'$ 对 $C$ 可容许. (使用\hyperref[I.prop8.8.2]{8.8.2}.)

\item 证明一个小的离散范畴对每个范畴 $C$ 都可容许.

\item 设 $J$ 是满足下列条件的范畴:
(i) $J$ 刚性, (ii) $\Ob J$ 可数, (iii) 对 $J$ 的任意两个对象 $j$, $j'$,
$\Hom(j,j')$ 有限, (iv) 对 $J$ 的每个对象 $j$, $J$ 中以 $j$ 为上界的对象的集合,
\ie 以 $j$ 为终点的某个箭头的起点的集合, 是有限的.
证明 $J$ 对每个 $𝒰$-范畴 $C$ 都可容许.
(先证明 $J$ 可以写成有限全子范畴 $J_n$ 的一个滤过并,
使得 $J$ 中以 $J_n$ 的某个对象为上界的每个对象都属于 $J_n$.
然后验证\hyperref[I.prop8.8.3]{8.8.3} 的判别条件 a), b).)
\end{enumerate}
\end{enonce*}

\begin{enonce*}[remark]{习题 8.8.7}\etiquette[8.8.7]{I.exer8.8.7}
在所用的记号中, 将有序集与它所定义的范畴等同.
\begin{enumerate}
\renewcommand{\theenumi}{\alph{enumi}}
\renewcommand{\labelenumi}{\theenumi)}
\item 设 $C$ 为有序集. 令 $C^{\widetilde{\phantom{,}}}$ 为 $C$ 中滤过的、
且满足 $x\in A$ 与 $y\leq x$ 蕴含 $y\in A$ 的子集 $A$ 所成的集合.
令 $L_\circ:C\rightarrow
C^{\widetilde{\phantom{,}}}$ 为如下映射:
给每个 $x\in
C$ 对应集合 $L_\circ(x)$, 它由满足 $y\leq x$ 的那些 $y\in A$ 组成.
证明 $L$ 是单射, 且 $C$ 上的序由 $C^{\widetilde{\phantom{,}}}$ 上的序诱导.
对每个 $A\in
C^{\widetilde{\phantom{,}}}$, 考虑包含
$f(A):A\rightarrow C$, %%
它是 $C$ 的一个 ind-对象; 对 $C$ 的每个 ind-对象 $\varphi:I\rightarrow
C$,
令 $g(\varphi)\in C^{\widetilde{\phantom{,}}}$ 为 $C$ 的如下子集: 由 $C$ 中以某个形如 $\varphi(i)$ 的元素为上界的元素组成.
证明\pageoriginale 由此得到两个互为拟逆的等价
\[
\xymatrix@C=1.5cm{ \Ind(C)
\ar@{<->}[r]^(.5){\displaystyle\mathop{f,g}_{\approx}}&
C^{\widetilde{\phantom{a}}}, }
\]
它们把函子 $L_\circ$ 变为典范函子 $L:C\rightarrow \Ind(C)$.

\item 设对每个 $x\in C$, 以 $x$ 为下界 (\resp 上界) 的元素集合有限.
证明 $C$ 的每个 ind-对象 (\resp pro-对象) 都本质常值,
而且对每个这样的 $\varphi:I\rightarrow C$ (\resp $\varphi:I^\circ\rightarrow C$),
存在 $i_\circ\in \Ob I$, 使得在 $I_{/i_\circ}$ 上诱导的函子为常值
(\ie 将每个箭头映为同构).

\item 设 $C\subset\NN^\circ\times
\NN$ 由满足 $i\geq
j$ 的自然数对 $(j,i)$ 组成.
证明 $\NN^{\widetilde{\phantom{,}}}$ 同构于从 $\NN$
(如 a) 中那样将其等同于 $\NN^{\widetilde{\phantom{,}}}$ 的一个有序子集)
添加一个最大元素 $\infty$ 后得到的有序集.
证明 $C^{\widetilde{\phantom{,}}}$ 同构于 $\NN^\circ\times
\NN^{\widetilde{\phantom{,}}}$.
考虑函子
$$
\varphi:\NN^\circ\longrightarrow C^{\widetilde{\phantom{,}}}\quav
\varphi(j)=(j,\infty){\quoi}.
$$
\end{enumerate}
\end{enonce*}

\noindent 证明:
\begin{enumerate}
\renewcommand{\labelenumi}{\theenumi)}
\item $\varphi$ 的投射极限在 $C^{\widetilde{\phantom{,}}}$ 中不可表,
尽管 $C$ 中的滤过投射极限可表 (而且由 b), 它们本质常值).

\item 对每个函子 $\psi:\NN^\circ\rightarrow
C$,
有 $\Hom(\psi,\varphi)=∅$, \textit{从而}函子 $\varphi$ 没有带指标表示.
\end{enumerate}

\begin{enonce*}[remark]{习题 8.8.8}\etiquette[8.8.8]{I.exer8.8.8}
设\pageoriginale $X=\NN\amalg \NN$ 为 $\NN$ 的两个副本的不交并,
$s$ 为 $X$ 的对称变换, 对每个 $i\in \NN$,
令 $X_i$ 为 $\NN$ 的子集 $\Delta_{i+1}\amalg\Delta_i$
(其中 $\Delta_i$ 表示 $\NN$ 中由满足 $j\leq
i$ 的 $j$ 组成的子集).
令 $C$ 为 $X$ 的子集范畴的如下子范畴: 对象为各 $X_i$,
箭头为由 $X$ 的恒等变换或其对称变换 $s$ 诱导的各 $X_i$ 之间的映射.
以同样方式定义范畴 $C^{\widetilde{\phantom{,}}}$, 但添加对象 $X$.

\begin{enumerate}
\renewcommand{\theenumi}{\alph{enumi}}
\renewcommand{\labelenumi}{\theenumi)}
\item 证明 $\Ind(C)$ 等价于 $C^{\widetilde{\phantom{,}}}$,
且典范函子 $C\rightarrow
\Ind(C)$ 对应于包含 $C\rightarrow
C^{\widetilde{\phantom{,}}}$.
证明不存在由 $C$ 的一个对象 $Y$ 与 $Y$ 的一个自同态 $f$ 组成的对 $(Y,f)$,
同时带有 $\Ind(C)$ 中从 $(Y,f)$ 到 $(X,s)$ 的一个带自同态对象的态射.

\item 推出: 若 $J$ 为只有一个对象、且除恒等箭头外仅有一个二阶箭头的范畴,
则由 $(X,s)$ 定义的函子 $J\rightarrow \Ind(C)$ 没有带指标表示.
\end{enumerate}
\end{enonce*}
 

% END INCLUDED SGA4_English_current_buildable_source_20260906/source/01/drafts/SGA4_I_LINES_4271_4625_English.texfrag


% BEGIN INCLUDED SGA4_English_current_buildable_source_20260906/source/01/drafts/SGA4_I_LINES_4626_4853_English.texfrag
\subsection{$\Ind(C)$ 的正合性质}\etiquette[8.9]{I.sec8.9}

\begin{enonce*}{命题 8.9.1}\etiquette[8.9.1]{I.prop8.9.1}
设 $C$ 为 $𝒰$-范畴.
\begin{enumerate}
\renewcommand{\theenumi}{\alph{enumi}}
\renewcommand{\labelenumi}{\theenumi)}
\item 典范函子 $L$ \hyperlink{sga4.exp.I.equation.8.2.4.8}{(8.2.4.8)} 和 $c$ \hyperlink{sga4.exp.I.equation.8.4.1}{(8.4.1)}
$$
C\xrightarrow{c}\Ind(C)\xrightarrow{L}\widehat{C}
$$
与投射极限可交换.

\item 假定 $C$ 中有限归纳极限可表, 且 $C$ 等价于一个小范畴.
则 $\Ind(C)$ 中小的投射极限可表.

\item 若\pageoriginale $C$ 中小的投射极限
(\resp 有限投射极限) 可表, 则 $\Ind(C)$ 中也如此.
设 $J$ 为有限且刚性的范畴, 或为离散范畴;
若 $C$ 中 $J$ 型投射极限可表, 则 $\Ind(C)$ 中同型的投射极限也可表.

\item $\Ind(C)$ 中小的滤过归纳极限 (由\hyperref[I.prop8.5.1]{8.5.1}, 它们可表)
是左正合的, \ie 与有限投射极限可交换.
\end{enumerate}
\end{enonce*}

\noindent \emph{证明}.
\begin{enumerate}
\renewcommand{\theenumi}{\alph{enumi}}
\renewcommand{\labelenumi}{\theenumi)}
\item $L$ 与投射极限可交换, 是因为它全忠实,
且 $\widehat{C}$ 中的投射极限是“逐个变元”计算的.
这意味着, 要验证 $\widehat{C}$ 中的态射投射系
$F\rightarrow F_i$ $(i\in I)$ 将 $F$ 表示为各 $F_i$ 的投射极限,
只须对每个 $X\in \Ob C$, 验证集合映射的投射系%%
$\Hom(X,F)\rightarrow \Hom(X,F_i)$ 满足类似断言.
但 $C\subset \Ind(C)$.

$c$ 与投射极限可交换, 形式地由 $L$ 与投射极限可交换,
以及复合 $L\circ
c:C\rightarrow \widehat{C}$ 也如此而得到.

\item 由 a), 所述断言等价于说 ind-可表预层的每个投射极限都是 ind-可表的.
这由判别准则\hyperref[I.them8.3.3]{8.3.3} (v) 立即得到,
因为左正合函子的投射极限左正合
(这又由“投射极限与投射极限可交换” (\hyperref[I.sec2.5.0]{2.5.0}) 得到).

\item 这\pageoriginale 形式地由 $\widehat{C}$ 的类似性质
(\hyperref[I.cor3.3]{3.3}) 得到, 因为 $L$ 保守 (它全忠实),
且与所考虑类型的极限可交换 (a) 及\hyperref[I.prop8.5.1]{8.5.1}).

\item 众所周知 (\cf \hyperref[I.prop2.3]{2.3}),
有限投射极限的可表性等价于下列类型的投射极限的可表性

\begin{center}
空 $, \ldots ,\quad \vcenter{ \xymatrix@C=.25cm@R=.5cm{
                     \bigdot \ar[rd]&       &\bigdot\ar[ld]\\
                             &\bigdot&
}}$
\end{center}

\noindent (对应于某些有限有序集),
而小的投射极限的可表性又等价于乘积及有限投射极限的可表性.
因此 (c) 中所述的第二个断言蕴含第一个断言.
\end{enumerate}

\emph{先设 $I$ 有限.} 使用关于将函子 $\varphi:J\rightarrow \Ind(C)$
表示为带指标形式%%
(\hyperref[I.sec8.5.4]{8.5.4}) 的结果\hyperref[I.prop8.8.5]{8.8.5}
\hypertarget{sga4.exp.I.equation.8.9.1.1}{}$$
\displaylines{\rlap{$(8.9.1.1)$}\hfill \phi:J\times I\longrightarrow
C{\quoi},\hfill}
$$
则 $\varprojlim \varphi$ 的存在性是下列更精确且更一般的结果的一个特例:

\begin{enonce*}{推论 \sisi{8.8.2}{8.9.2}}\etiquette[8.9.2]{I.cor8.9.2}
考虑由带指标形式 $\Phi$ \hyperlink{sga4.exp.I.equation.8.9.1.1}{(8.9.1.1)}
给出的函子 $\varphi:J\rightarrow \Ind(C)$, 其中 $J$ 为有限范畴.
若对每个 $i\in \Ob I$, 偏函子 $j\rightarrow
\Phi(j,i)$ 的投射
(\resp 归纳) 极限在 $C$ 中可表, 则 $\varphi$ 的投射
(\resp 归纳) 极限在 $\Ind(C)$ 中可表, 且有典范同构
\hypertarget{sga4.exp.I.equation.8.9.2.1}{}$$
\displaylines{\rlap{$(8.9.2.1)$}\hfill
\varprojlim\varphi\simeq\textrm{« }{\varinjlim_i}\textrm{ »} \varprojlim_j
\Phi(j,i)\hfill}
$$
(\resp\pageoriginale
\hypertarget{sga4.exp.I.equation.8.9.2.2}{}$$
\displaylines{\rlap{$(8.9.2.2)$} \hfill \varinjlim \varphi \simeq
\textrm{« }{\varinjlim_i}\textrm{ »} \varprojlim_j\Phi(j,i)){\quoi},\hfill}
$$
其中 $\varprojlim_j$ (\resp. $\varinjlim_j$) 在 $C$ 中计算.
\end{enonce*}

对第一个公式, 只须利用 $\Ind(C)$ 中滤过归纳极限与 $\varprojlim_j$ 可交换,
以及 $c$ 与 $\varprojlim_J$ 可交换
\hyperref[I.prop8.9.1]{(8.9.1 d)} 和 a)):
$$
\varprojlim \varphi=\varprojlim_j \textrm{« }{\varinjlim_i}\textrm{ »\,} \Phi
(j,i) \xleftarrow{\sim} \textrm{« }{\varinjlim_i}\textrm{ »\,}
{\varprojlim_j}^{\Ind(C)} \Phi(j,i)\simeq \textrm{« }{\varinjlim_i}\textrm{ »}%%
\varprojlim_j{^C \ \Phi(j,i){\quoi}}.
$$
第二个公式同样证明, 使用 $I$ 型归纳极限与 $J$ 型归纳极限可交换
(\hyperref[I.sec2.5.0]{2.5.0}), 以及 $c$ 与 $J$ 型归纳极限可交换这一事实;
后者在下面的\hyperref[I.rem8.9.4]{8.9.4} a) 中证明.

\emph{$J$ 离散的情形}. 这一情形包含在下列更精确的断言中:

\begin{enonce*}{推论 8.9.3}\etiquette[8.9.3]{I.cor8.9.3}
设各 $I_\alpha$ 为以小集合 $A$ 为指标的本质小的滤过范畴,
且对每个 $\alpha \in A$, 设 $X(\alpha)=(X(\alpha)_{i_\alpha})_{i_\alpha \in I_\alpha}$
为 $C$ 的以 $I_\alpha$ 为指标的 ind-对象.
令 $I$ 为各 $I_\alpha$ 的乘积范畴,
并设对每个 $\sheaf{i}=(i_\alpha)_{\alpha\in A}\in I$,
乘积 $\prod_{\alpha\in A}%%
X (\alpha)_{i_\alpha}$ 在 $C$ 中可表.
则乘积 $\prod_{\alpha \in A} %%
X (\alpha)$ 在 $\Ind(C)$ 中可表,
并典范同构于由下式给出的以 $I$ 为指标的 ind-对象
\hypertarget{sga4.exp.I.equation.8.9.3.1}{}$$
\displaylines{\rlap{$(8.9.3.1)$}\hfill \prod_{\alpha \in A}\;
\mathop{\textrm{« }\varinjlim\textrm{ »}}_{i_\alpha\in I_\alpha}\;
X(\alpha)_{i_\alpha}\simeq
\mathop{\textrm{« }\varinjlim\textrm{ »}}_{\sheaf{i}=(i_\alpha)_{\alpha\in
A}\in I}\; \left(\prod_{\alpha\in
A}A(\alpha)_{i_\alpha}\right){\quoi},\hfill}
$$
其中\pageoriginale 右边的乘积是在 $C$ 中计算的乘积.
(注意 $I$ 滤过且本质小, 因为各 $I_\alpha$ 如此.)
\end{enonce*}

事实上, 众所周知, 当各极限在 $\widehat{C}$ 中计算时, 上式成立
(将其归结到集合范畴的情形, 即可立即看出).
于是由 $L$ 与所考虑的极限可交换
(\hyperref[I.prop8.9.1]{8.9.1} a) 和\hyperref[I.prop8.5.1]{8.5.1}) 得到结论.

\begin{enonce*}[remark]{注 8.9.4}\etiquette[8.9.4]{I.rem8.9.4}
\hyperref[I.cor8.9.2]{8.9.2} 和\hyperref[I.cor8.9.3]{8.9.3} 的证明更一般地表明,
若对每个 $i\in \Ob I$, \emph{在 $\Ind(C)$ 中计算的}
$\varprojlim_j \Phi(j,i)$ (\resp $\varinjlim_j \Phi(j,i)$) 可表,
则 $\varprojlim_i\varphi$ (\resp $\varinjlim_j\varphi$) 也如此.
将这一点与 c) 的论证结合可知: 对给定的范畴%%
$J$, 若它来自一个有限或离散有序集 (或更一般地, 它对 $C$ 可容许
(\hyperref[I.prop8.3.1]{8.3.1})), 则 $\Ind(C)$ 中 $J$ 型投射极限
(\resp 归纳极限) 可表, 只要 (而且仅当) 对每个函子 $\varphi:J\rightarrow C$,
$\varphi$ 的\emph{在 $\Ind(C)$ 中计算的}投射 (\resp 归纳) 极限可表.
在不带括号的情形, 由 (a), 这也就是说 $C$ 上可表预层的每个 $J$ 型投射极限均 ind-可表.
\end{enonce*}

\begin{enonce*}{命题 8.9.5}\etiquette[8.9.5]{I.prop8.9.5}
设 $C$ 为 $𝒰$-范畴.
\begin{enumerate}
\renewcommand{\theenumi}{\alph{enumi}}
\renewcommand{\labelenumi}{\theenumi)}
\item 典范函子
$$
c:C\longrightarrow \Ind(C)
$$
右正合 (因而由 $\hyperref[I.prop8.9.1]{8.9.1}$ {\rm a)}, 它正合).

\item 若 $C$ 中有限归纳极限 (\resp 有限和) 可表,
则 $\Ind(C)$ 中小的归纳极限 (\resp 小的和) 可表.
设 $J$ 为有限或离散的预序\pageoriginale 集;
若 $C$ 中 $J$ 型归纳极限可表, 则 $\Ind(C)$ 中也如此.
\end{enumerate}
\end{enonce*}

\noindent \emph{证明.}
\begin{enumerate}
\renewcommand{\theenumi}{\alph{enumi}}
\renewcommand{\labelenumi}{\theenumi)}
\item 设在 $C$ 中有 $X\simeq \varinjlim_J X_j$,
其中 $J$ 为有限范畴, $X$ 及各 $X_j$ 属于 $C$.
则对 $C$ 的每个 ind-对象 $𝒴=(Y_i)_{i\in I}$, 有:
\begin{gather*}
\Hom(X,𝒴)\simeq \varinjlim_i \Hom(X,Y_i)\simeq \varinjlim_i
\varprojlim_j \Hom(X_j,Y_i)\\
\mathrel{\mathop\simeq^{\hyperref[I.prop2.8]{2.8}}}\varprojlim_j \varinjlim_i \Hom(X_j,
Y_i) \simeq \varprojlim_j \Hom (X_j,𝒴){\quoi},
\end{gather*}
比较两端即得所要的结论.

\item 我们已在\hyperref[I.rem8.8.4]{8.8.4} 中指出,
至少对有限归纳极限而言, 所述断言由 a) 和\hyperref[I.cor8.9.2]{8.9.2} 得到.
要在无限情形下证明所述结论, 可归结到和的情形;
和可解释为有限和的一个滤过归纳极限, 因而由\hyperref[I.prop8.5.1]{8.5.1} 得出结论.
\end{enumerate}

% END INCLUDED SGA4_English_current_buildable_source_20260906/source/01/drafts/SGA4_I_LINES_4626_4853_English.texfrag


% BEGIN INCLUDED SGA4_English_current_buildable_source_20260906/source/01/drafts/SGA4_I_LINES_4854_5078_English.texfrag
\begin{enonce*}[remark]{注 8.9.6}\etiquette[8.9.6]{I.rem8.9.6}
我们已指出, 函子 $c$ 一般不与滤过归纳极限可交换,
因而也不与无限和可交换; 这与 $L:\Ind(C)\rightarrow
\widehat{C}$ 的情形相反.
另一方面, 除了与\emph{滤过}归纳极限可交换 (\hyperref[I.prop8.5.1]{8.5.1}) 外,
$L$ 与任何其他类型的归纳极限几乎都没有可交换性
(始对象、两个对象的和、双箭头的余核).
因此, 若 $C$ 有始对象 $∅_C$, 由 a), 它也是 $\Ind(C)$ 的始对象;
$∅_C$ 却绝不是 $\widehat{C}$ 的始对象
(\ie 等同于值为 $∅$ 的常值预层), 因为 $\Hom(∅_C,∅_C)\neq ∅$!
\end{enonce*}

\begin{enonce*}{命题 8.9.7}\etiquette[8.9.7]{I.prop8.9.7}
设\pageoriginale
$$
f:C\longrightarrow C'
$$
为 $𝒰$-范畴之间的函子, $J$ 为有限刚性范畴 $\hyperref[I.prop8.8.5]{(8.8.5)}$.
若 $C$ 中 $J$ 型归纳 (\resp 投射) 极限可表, 且 $f$ 与它们可交换,
则 $\Ind(f):\Ind(C)\rightarrow \Ind(C')$ 也与这一类型的极限可交换.
\end{enonce*}

这由\hyperref[I.prop8.8.5]{8.8.5} 以及\hyperref[I.cor8.9.2]{8.9.2}
关于带指标形式表示的函子在 ind-对象范畴中的有限极限的计算立即得到.

\begin{enonce*}{推论 8.9.8}\etiquette[8.9.8]{I.cor8.9.8}
若 $C$ 中有限归纳 (\resp 投射) 极限可表, 且 $f$ 右正合 (\resp 左正合),
则 $\Ind(f)$ 也如此; 在不带括号的情形, $\Ind(f)$ 甚至与所有小的归纳极限可交换.
\end{enonce*}

第一个断言由\hyperref[I.prop8.9.7]{8.9.7} 得到.
考虑到 $\Ind(f)$ 与滤过归纳极限可交换 (\hyperref[I.prop8.6.3]{8.6.3}),
第二个断言由第一个断言得到.

\begin{enonce*}[remark]{习题 8.9.9}\etiquette[8.9.9]{I.exer8.9.9}
设 $C$ 为 $𝒰$-范畴.
\begin{enumerate}
\renewcommand{\theenumi}{\alph{enumi}}
\renewcommand{\labelenumi}{\theenumi)}
\item 设 $C$ 中有限和可表.
证明: 若 $C$ 中\emph{有限}和不交 (\resp 普遍)
(\cf II \hyperref[II.def4.5]{4.5}),
则 $\Ind(C)$ 中\emph{小的}和不交 (\resp 普遍).

\item 设 $C$ 中每个态射都可分解为一个满态射后接一个单态射
(\resp 一个有效满态射后接\pageoriginale 一个单态射,
\resp 一个满态射后接一个有效单态射,
\resp 一个有效满态射后接一个有效单态射).
证明 $\Ind(C)$ 满足同样的性质.
在最后一种情形下, 由此可知 $\Ind(C)$ 的每个满态射有效,
$\Ind(C)$ 的每个单态射有效, 且 $\Ind(C)$ 的每个双态射都是同构;
证明在这一情形下, $\Ind(C)$ 的每个满态射 (\resp 单态射)
可表示为 $C$ 中满态射 (\resp 单态射) 的一个归纳系.
若进一步假定 $C$ 中每个满态射 (\resp 每个单态射) 都\emph{普遍},
则 $\Ind(C)$ 中也有同样的性质.

\item 假定 $C$ 是加性 (\resp 阿贝尔) 的; 则 $\Ind(C)$ 也如此.

\item 设 $X=\textrm{« }\varinjlim_I\textrm{ »\,}X_i$ 为 $C$ 的一个 ind-对象,
$R\rightrightarrows X$ 为 $X$ 上的一个等价关系.
对每个 $i\in \Ob I$, 令 $R_i\rightrightarrows X_i$ 为经 $X_i\rightarrow X$
由 $R$ 诱导的 $X_i$ 上的等价关系 (将其看作 $\Ind(C)$ 的对象).
(我们假定纤维积 $R_i=R\times_{X_XX}X_i \times X_i$ 在 $(\Ind(C))$ 中可由
$\Ind(C)$ 的一个对象表示; 当 $C$ 中有限投射极限可表时便如此.)
证明: 要使 $R$ 有效, 只须各 $R_i$ 有效.

\item 设 $X$ 是 $C$ 的一个对象, $R$ 为 $c(X)$ 上的一个等价关系,
\ie 为视作 $\Ind(C)$ 的对象的 $X$ 上的等价关系.
假定 $C$ 中纤维积可表. 要使 $R$ 有效, 必须使 $R$ 具有
$\textrm{« }\varinjlim_i\textrm{ »\,} R_i$ 的形式,
其中各 $R_i$ 为 $X$ (视作 $C$ 的对象) 上的等价关系;
若假定 $C$ 中等价关系有效, 则这个条件\pageoriginale 也是充分的.

\item 假定 $C$ 中每个态射都可分解为一个有效满态射后接一个有效单态射,
且有限投射极限可表. 设 $X$ 是 $C$ 的对象,
$R$ 为 $X\times X$ 的子对象, 视作 $\Ind(C)$ 的对象.
将 $R$ 写成 $\textrm{« }\varinjlim\mbox{''
}Z_i$ 的形式, 其中 $(Z_i)_{i\in I}$ 为 $C$ \emph{中}
$X\times X$ 的子对象的滤过递增族 (由 b), 这可以做到).
证明: $R$ 为 $X$ 上的等价关系的充要条件是,
对每个 $i\in \Ob I$, 存在指标 $j\in \Ob I$,
其对应子对象包含 $s(Z_i)$ 与 $Z_i\circ Z_i$,
其中 $s$ 为 $X\times X$ 的对称变换.

\item 假定 $C$ 中有限投射极限及有限和可表, 每个等价关系有效,
且每个态射都可分解为一个有效满态射后接一个有效单态射.
证明 $\Ind(C)$ 中等价关系普遍, 当且仅当 $C$ 满足下列条件:
\begin{itemize}
\item[%%
ST)] 对 $C$ 的每个对象 $X$ 和 $C$ 中 $X\times
X$ 的每个子对象 $Z$,
若按 $Z_\circ=Z$,
$Z_n=\Sup(Z_n,s(Z_n),Z_n\circ Z_n)$ 递归定义 $C$ 中 $X\times X$ 的子对象序列
$Z_n$ $(n\geq 0)$
(其中 $\Sup$ 在 $X\times X$ 的子对象集合中取, 由对 C 的假设它存在),
则序列 $(Z_n)_{n\geq 0}$ \emph{稳定}.
\end{itemize}

\setcounter{enumi}{10}
\item 证明 $\Ind((\Ens))$ 中的等价关系不一定有效.

$^\ast$ NB. 关于保证 $\Ind(C)$ 为拓扑斯的判别准则, \cf VI
8.9.9.${}_\ast$
\end{enumerate}
\end{enonce*}

\begin{enonce*}[remark]{习题 8.9.10}\etiquette[8.9.10]{I.exer8.9.10}
设\pageoriginale $C$ 为 $𝒰$-范畴. 置
$\Pro(C)=(\Ind(C^\circ))^\circ$ (\cf \hyperref[I.sec8.11]{8.11}).
\begin{enumerate}
\renewcommand{\theenumi}{\alph{enumi}}
\renewcommand{\labelenumi}{\theenumi)}
\item 假定 $C$ 中有限和 (\resp 小的和) 可表.
证明: 若它们不交 (\cf II \hyperref[II.def4.5]{4.5}),
则 $\Pro(C)$ 满足同样的%%
条件.

\item 设 $J$ 为小集合, 使得 $C$ 中以 $J$ 为指标的和可表,
因而 $\Pro(C)$ 中也如此.
设 $(X(\alpha))_{\alpha\in J}$ 为 $\Pro(C)$ 的对象族, 其中
$X(\alpha)=\textrm{« }\varprojlim_{I_\alpha}\textrm{ »\,}X_{i_\alpha}$.
证明: 为使 $\Ind(C)$ 中各 $X(\alpha)$ 的和普遍,
只须每族 $\varprojlim_J X_{i_\alpha}$ 都如此,
其中 $(i_\alpha)_{a\in I}\in
\prod_{\alpha\in J}I_\alpha$.
推出: 若 $C$ 中 $J$ 型的和普遍, 则 $\Pro(C)$ 中也如此的充要条件是,
对上述每个族 $(X(\alpha))_{\alpha\in J}$,
若对每个 $\alpha\in J$ 有 $I_\alpha=I$,
则 $\Pro(C)$ 中的典范同态
$$
\textrm{« }\varprojlim_I\textrm{ »\,}\amalg_\alpha X(\alpha)_i\longleftarrow
\textrm{« }\varprojlim_{I^J}\textrm{ »\,} \amalg_{\alpha\in
J}X(\alpha)_{i_\alpha}
$$
为同构. 推出 $\Pro((\Ens))$ 中 $J$ 型的和普遍, 当且仅当 $J$ 有限.
\end{enumerate}
\end{enonce*}

\begin{enonce*}[remark]{习题 8.9.11}\etiquette[8.9.11]{I.exer8.9.11}
设 $C$ 为 $𝒰$-范畴.
\begin{enumerate}
\renewcommand{\theenumi}{\alph{enumi}}
\renewcommand{\labelenumi}{\theenumi)}
\item 设 $(X_i\rightarrow X)_{i\in I}$ 为 $C$ 中的态射族.
证明: 要使它在 $\Ind(C)$ 中为满态射族, 只须存在一个在 $C$ 中为满态射族的\emph{有限}子族.
证明: 当假定 $C$ 中以 $X_i$ 为终点的每个有限态射族都可分解为一个满态射族\pageoriginale
后接一个有效单态射\hyperref[I.sec10.5]{(10.5)} 时, 这个条件也是必要的.
(对最后这个断言, 令 $P$ 为 $I$ 的有限子集按包含关系排序所成的集合,
令 $X'=(X'_J)_{J\in P}$ 为给定族的有限子族的像所组成的 ind-对象;
最后令 $X''=X\amalg_{X'}X=\textrm{« }\lim\textrm{ »\,}X\amalg_{X'_J}X$.
证明两个典范态射 $X\rightrightarrows X''$ 不同, 但在各 $X'_J$ 上相同.)

\item 假定 $C$ 中具有同一终点的每个有限态射族都可分解为一个满态射族后接一个有效单态射.
证明: $\Ind(C)$ 有一个通过满态射生成的小子范畴 (\hyperref[I.def7.1]{7.1}),
当且仅当 $C$ 有小子范畴 $C'$, 使得 $C$ 的每个对象都是一个起点在 $C'$ 中的\emph{有限}满态射族的终点.
若假定 $C$ 有小的生成子范畴 (\hyperref[I.def7.1]{7.1}),
则 $\Ind(C)$ 有通过满态射生成的小子范畴, 当且仅当 $C$ 等价于一个小范畴.
(对最后这个断言, 使用\hyperref[I.var7.5.2]{7.5.2}).

\item $\Ind((\Ens))$ 没有小的生成子范畴.
\end{enumerate}
\end{enonce*}

\begin{enonce*}[remark]{习题 8.9.12}\etiquette[8.9.12]{I.8.9.12}
设 $C$ 为 $𝒰$-范畴, 考虑
$\Pro(C)=\Ind(C^\circ)^\circ$.
\begin{enumerate}
\renewcommand{\theenumi}{\alph{enumi}}
\renewcommand{\labelenumi}{\theenumi)}
\item 证明: 若 $\Pro(C)$ 有一个通过满态射生成的小子范畴 $P'$
\hyperref[I.def7.1]{(7.1)}, 则 $C$ 有一个小子范畴 $C'$,
它在 $\Pro(C)$ 中通过满态射生成.
(取 $P'$ 的各对象的分量所成的范畴.)

\item 证明\pageoriginale 范畴 $\Pro((\Ens))$ 没有通过满态射生成的小子范畴.
(若 $C'$ 如 a) 中那样, 选取一个集合 $X$, 其基数严格大于 $C'$ 的各对象的基数,
并考虑由 $X$ 中基数 $<\card(X)$ 的子集的补集组成的 pro-集合 $𝒳$.
证明: 对基数 $<\card(X)$ 的每个非空集合 $Y$, 有 $\Hom(Y,𝒳)=∅$,
但 $𝒳$ 不同构于常值 pro-集合 $∅$.)
\end{enumerate}
\end{enonce*}

% END INCLUDED SGA4_English_current_buildable_source_20260906/source/01/drafts/SGA4_I_LINES_4854_5078_English.texfrag


% BEGIN INCLUDED SGA4_English_current_buildable_source_20260906/source/01/drafts/SGA4_I_LINES_5079_5229_English.texfrag
\subsection{对偶概念: pro-对象、pro-可表函子}\etiquette[8.10]{I.sec8.10}

设 $C$ 为 $𝒰$-范畴. $C$ 的一个 \emph{pro-对象}是一个函子
\hypertarget{sga4.exp.I.equation.8.10.1}{}$$
\displaylines{\rlap{$(8.10.1)$}\hfill \varphi:I^\circ\longrightarrow
C{\quoi},\hfill}
$$
其中 $I$ 是本质小的滤过范畴 (称为指标范畴),
如常, $I^\circ$ 表示其对偶范畴.
设 $𝒱$ 为满足 $𝒱\supset 𝒰$ 的宇宙.
注意到 $C$ 的以 $I\in 𝒱$ 为指标的 pro-对象,
与 $C^\circ$ 的以 $I\in
𝒱$ 为指标的 ind-对象一一对应:
给每个这样的 ind-对象 $\psi:I\rightarrow C^\circ$
对应 pro-对象 $\varphi=\psi^\circ:I^\circ\rightarrow C$.
沿着这一对应, 从 ind-对象的类似概念
\hyperlink{sga4.exp.I.equation.8.2.4.2}{(8.2.4.2)},
\hyperlink{sga4.exp.I.equation.8.2.5.1}{(8.2.5.1)} 出发,
“通过结构的搬移和箭头的反向”定义 $C$ 的 pro-对象之间的态射概念,
从而定义 $C$ 的以 $I\in 𝒱$ 为指标的 pro-对象范畴, 并有典范同构
\hypertarget{sga4.exp.I.equation.8.10.2}{}$$
\displaylines{\rlap{$(8.10.2)$}\hfill
\Pro_{𝒱}(C,𝒰)\simeq \Ind_{𝒱}(C^\circ,
𝒰)^\circ{\quoi};\hfill}
$$
当\pageoriginale $𝒱=𝒰$ 时, 就称为 \emph{$C$ 的 pro-对象范畴},
记为 $\Pro_{𝒰}(C)$ 或简记为 $\Pro(C)$,
因此由下式定义:
\hypertarget{sga4.exp.I.equation.8.10.3}{}$$
\displaylines{\rlap{$(8.10.3)$}\hfill
\Pro(C)=\Pro_{𝒰}(C,𝒰)\simeq
\Ind(C^\circ)^\circ{\quoi}.\hfill}
$$
若两个 pro-对象以\emph{带指标形式}给出
\hypertarget{sga4.exp.I.equation.8.10.4}{}\begin{equation}
𝒳=(X_i)_{i\in I}\quav 𝒴=(Y_j)_{j\in
J}{\quoi},\tag{8.10.4}
\end{equation}
则公式\hyperlink{sga4.exp.I.equation.8.2.5.1}{(8.2.5.1)} 在这里变为
\hypertarget{sga4.exp.I.equation.8.10.5}{}\begin{equation}
\Hom(𝒳,𝒴)\simeq \varprojlim_j \varinjlim_i
\Hom(X_i,Y_j){\quoi}.\tag{8.10.5}
\end{equation}
将函子\hyperlink{sga4.exp.I.equation.8.4.1}{(8.4.1)} 应用于 $C^\circ$,
再取对偶范畴, 得到一个典范函子 (无混淆之虞时仍用同一字母 $c$ 表示),
它将 $C$ 等同于 $\Pro(C)$ 的一个全子范畴:
\hypertarget{sga4.exp.I.equation.8.10.6}{}\begin{equation}
c:C\longrightarrow\Pro(C){\quoi};\tag{8.10.6}
\end{equation}
正是为了有这样的函子, 而不是 $C\rightarrow\Pro(C)^\circ$,
我们才在从 ind-对象的类似定义出发定义 pro-对象的态射时“反转箭头”.
\hyperlink{sga4.exp.I.equation.8.10.6}{(8.10.6)} 的本质像中的 pro-对象仍称为\emph{本质常值 pro-对象}.

置
\hypertarget{sga4.exp.I.equation.8.10.7}{}\begin{equation}
\check{C}=(C^\circ)^{\widehat{\phantom{,}}}=\Hhom(C,(\Ens)){\quoi},\tag{8.10.7}
\end{equation}
则可将函子\hyperlink{sga4.exp.I.equation.8.2.4.7}{(8.2.4.7)} 看作典范的全忠实函子
\hypertarget{sga4.exp.I.equation.8.10.8}{}\begin{equation}
L:\Pro_{𝒱}(C,𝒰)\longrightarrow
\check{C}^\circ{\quoi},\tag{8.10.8}
\end{equation}\pageoriginale
特别地得到
\hypertarget{sga4.exp.I.equation.8.10.9}{}\begin{equation}
L:\Pro(C)\hookrightarrow \check{C}^\circ{\quoi}.\tag{8.10.9}
\end{equation}

我们把按需要将前面及后面各节的结果从 ind-对象的语言翻成 pro-对象的语言的工作留给读者.
根据具体的数学背景, 可能这一种语言更有用, 也可能另一种更有用,
更不用说两种概念同时出现的情形, 例如必须考虑
$\Pro(\Ind(C))$ 或 $\Ind(\Pro(C))$ 这样的复杂范畴时.
这里只再作几点说明, 以固定术语与记号, 并将有关“演算”说清楚.


\setcounter{subsubsection}{9}
\subsubsection{}\etiquette[8.10.10]{I.sec8.10.10}
一个\emph{协变函子 $F:C\rightarrow (\Ens)$}, 若属于
\hyperlink{sga4.exp.I.equation.8.10.9}{(8.10.9)}
(或等价地, \hyperlink{sga4.exp.I.equation.8.10.8}{8.10.8}) 的本质像,
则称为 \emph{pro-可表函子}; 因而由这些函子组成的 $\check{C}$ 的全子范畴等价于 $\Pro(C)^\circ$.
一般更愿意在对偶范畴中工作, 即在\hyperlink{sga4.exp.I.equation.8.10.9}{(8.10.9)}
\emph{作为子范畴的}本质像中工作, 因而它等价于 $\Pro(C)$ 本身.
遵循这一用法, 常常以将协变函子
$$
F:C\longrightarrow (\Ens)
$$
看作 $\Hhom(C,(\Ens))^\circ=\check{C}^\circ$ 的对象为宜,
从而将态射
$$
X\longrightarrow F\mbox{ 在 } \check{C}^\circ{\quoi},
$$
的范畴, \ie\pageoriginale 由 $C$ 的一个对象 $X$ 与一个元素 $u\in F(X)$
组成的对 $(X,u)$ 的范畴, 记为
\hypertarget{sga4.exp.I.equation.8.10.11}{}\begin{equation}
{}_{\displaystyle F}\backslash C{\quoi}  %%
\ .\tag{8.10.11}
\end{equation}
在这一约定下, 于是得到一个\emph{协变}函子 (起点函子)
\hypertarget{sga4.exp.I.equation.8.10.12}{}\begin{equation}
{}_{\displaystyle F}\backslash C\rightarrow C{\quoi},\tag{8.10.12}
\end{equation}
而\hyperref[I.prop3.4]{3.4} 改写为
\hypertarget{sga4.exp.I.equation.8.10.13}{}\begin{equation}
F\xrightarrow{\sim}%%
\textrm{``}\varprojlim_{({}_{\displaystyle F}\backslash C)^\circ}\textrm{''\,}
X\tag{8.10.13}
\end{equation}

\setcounter{subsubsection}{13}
\subsubsection{}\etiquette[8.10.14]{I.sec8.10.14}
将判别准则\hyperref[I.them8.3.3]{8.3.3} 应用于 $C^\circ$,
得到 $F$ 的 pro-可表性判别准则:
充要条件是 $({}_F\backslash C)^\circ$ 滤过且本质小;
若 $C$ 等价于一个小范畴, 则后一个条件多余;
若进一步 $C$ 中有限投射极限可表,
则这又等价于 $F$ 与它们可交换, \ie \emph{左正合}.

\subsubsection{}\etiquette[8.10.15]{I.sec8.10.15}
在 $\Pro(C)$ 中, 小的滤过投射极限可表,
且函子\hyperlink{sga4.exp.I.equation.8.10.9}{(8.10.9)} 与它们可交换;
换句话说, 这个函子将 $\Pro(C)$ 中的投射极限映到 $\check{C}=\Hhom(C,(\Ens))$
中的\emph{归纳}极限, 就像它与\hyperlink{sga4.exp.I.equation.8.10.6}{(8.10.6)}
的复合 $C\rightarrow\check{C}^\circ$ 一样;
两者都具有如下令人不快的性质: 看作取值于 $\check{C}$ 时是反变的.
应注意, 从 $C$ 到 $(\Ens)$ 的 pro-可表\pageoriginale 函子,
是可表函子的小的滤过\emph{归纳}极限
(并非投射极限, 尽管术语可能使人这样想).


% END INCLUDED SGA4_English_current_buildable_source_20260906/source/01/drafts/SGA4_I_LINES_5079_5229_English.texfrag


% BEGIN INCLUDED SGA4_English_current_buildable_source_20260906/source/01/drafts/SGA4_I_LINES_5230_5389_English.texfrag
\subsection{ind-伴随与 pro-伴随}\etiquette[8.11]{I.sec8.11}

\subsubsection{}\etiquette[8.11.1]{I.sec8.11.1}
设
\hypertarget{sga4.exp.I.equation.8.11.1.1}{}\begin{equation}
f:C\longrightarrow C'\tag{8.11.1.1}
\end{equation}
是 $𝒰$-范畴之间的函子, 从而有一个函子
$F \sisi{→}{↦} F\circ f$
\hypertarget{sga4.exp.I.equation.8.11.1.2}{}\begin{equation}
f^\ast:\widehat{C}'\longrightarrow\widehat{C}{\quoi}.\tag{8.11.1.2}
\end{equation}
若上面的函子将 ind-可表函子映为 ind-可表函子,
则称 $f$ \emph{有一个 ind-伴随}.
由于 $f^\ast$ 与归纳极限可交换,
且由 ind-可表函子组成的 $\widehat{C}'$ 的全子范畴对小的滤过归纳极限稳定,
说 $f$ 有 ind-伴随, 等价于说 $f^\ast$ 将 $C'$ 的对象
(将其等同于 $\widehat{C}'$ 的一个全子范畴) 映为 $C$ 上的 ind-可表函子, \ie
\hypertarget{sga4.exp.I.equation.8.11.1.3}{}\begin{equation}
X\longmapsto \Hom(f(X),Y')\mbox{ 对每个 }
Y'\in \Ob C'\mbox{ 都 ind-可表}{\quoi}.\tag{8.11.1.3}
\end{equation}
还可以用另一种方式表述 $f$ 有 ind-伴随的条件:
存在一个本质上由 $f^\ast$ 诱导的函子
\hypertarget{sga4.exp.I.equation.8.11.1.4}{}\begin{equation}
g:\Ind(C')\longrightarrow \Ind(C)\tag{8.11.1.4}
\end{equation}
\ie 存在一个双函子同构
\hypertarget{sga4.exp.I.equation.8.11.1.5}{}\begin{equation}
\Hom_{\Ind(C)}(X,g(𝒴'))\simeq
\Hom_{\Ind(C')}(f(X),𝒴'){\quoi},\tag{8.11.1.5}
\end{equation}
其中\pageoriginale $X\in \Ob C$ 而 $𝒴'\in \Ob \Ind(C')$.
事实上, 甚至只须有一个函子
\hypertarget{sga4.exp.I.equation.8.11.1.6}{}\begin{equation}
g_\circ:C'\longrightarrow \Ind(C)\tag{8.11.1.6}
\end{equation}
以及一个关于 $X\in \Ob C$, $Y'\in \Ob C'$ 的双函子同构
\hypertarget{sga4.exp.I.equation.8.11.1.7}{}\begin{equation}
\Hom_{\Ind(C)}(X,g_\circ(Y'))\simeq \Hom_{C'}(f(X),Y')\tag{8.11.1.7}
\end{equation}
即可. 显然, 函子 $g$ (分别地, $g_\circ$) 由 $f$ 决定到唯一同构,
反过来由这个 $g$ (分别地, $g_\circ$) 又能将 $f$ 重构到典范同构.
由\hyperlink{sga4.exp.I.equation.8.11.1.5}{(8.11.1.5)} 可见,
函子 $g$ 与归纳极限可交换, 并且它“延拓”函子 $g_\circ$;
因此 $g_\circ$ 将它决定到唯一同构 (\hyperref[I.sec8.7.2]{8.7.2}).
函子 $g$, 有时也包括它所延拓的函子 $g_\circ$,
称为 \emph{$f$ 的 ind-伴随函子}
(这里不必说明“右伴随”, 因为若另一个伴随存在, 它将称为 pro-伴随;
见下面的\hyperref[I.sec8.11.5]{8.11.5}).

当然, 若 $f$ 有右伴随
$$
f^{\ad}:C'\longrightarrow C{\quoi},
$$
则它有 ind-伴随 $g$, 且后者典范同构于 $f^{\ad}$ 在 ind-对象上的典范延拓:
\hypertarget{sga4.exp.I.equation.8.11.1.8}{}\begin{equation}
g\simeq\Ind(f^{\ad}){\quoi}.\tag{8.11.1.8}
\end{equation}
因此 ind-伴随的概念是右伴随概念的自然推广.

现在考虑\pageoriginale 典范延拓
\hypertarget{sga4.exp.I.equation.8.11.1.9}{}\begin{equation}
\Ind(f):\Ind(C)\longrightarrow\Ind(C'){\quoi}.\tag{8.11.1.9}
\end{equation}
于是有:

\begin{enonce*}{命题 8.11.2}\etiquette[8.11.2]{I.prop8.11.2}
函子 $f:C\rightarrow C'$ 有 ind-伴随的充要条件是函子
$\Ind(f)$ $\hyperlink{sga4.exp.I.equation.8.11.1.9}{(8.11.1.9)}$ 有右伴随.
后者典范同构于 ind-伴随 $\hyperlink{sga4.exp.I.equation.8.11.1.4}{(8.11.1.4)}$.
\end{enonce*}

充分性及最后一个断言立即由 ind-伴随公式
\hyperlink{sga4.exp.I.equation.8.11.1.5}{(8.11.1.5)} 得到.
关于必要性, 注意到对各 $X_i$ 取这一公式的投射极限,
便得到关于 ind-对象 $𝒳=(X_i)_{i\in I}$ 和 $𝒴'$ 的伴随同构:
\hypertarget{sga4.exp.I.equation.8.11.2.1}{}\begin{equation}
\Hom_{\Ind(C)}(𝒳, g(𝒴'))\simeq
\Hom_{\Ind(C')}(\Ind(f)(𝒳),𝒴'){\quoi},\tag{8.11.2.1}
\end{equation}
\cqfd

\begin{enonce*}{推论 8.11.3}\etiquette[8.11.3]{I.cor8.11.3}
若 $f$ 有 ind-伴随, 则 $\Ind(f)$ 与归纳极限可交换,
而 ind-伴随 $g$ 与投射极限可交换.
\end{enonce*}

\begin{enonce*}{命题 8.11.4}\etiquette[8.11.4]{I.prop8.11.4}
函子 $f:C\rightarrow C'$ 有 ind-伴随的必要条件是 $f$ \sisi{左}{右}正合,
而当 $C$ 等价于一个小范畴时, 此条件也是充分的.
\end{enonce*}

这由判别准则\hyperlink{sga4.exp.I.equation.8.11.1.3}{(8.11.1.3)},
以及\hyperref[I.prop8.3.1]{8.3.1} 和\hyperref[I.them8.3.3]{8.3.3} (iv) 得到.

关于用可及函子的概念表述的另一个判别准则,
见\hyperref[I.cor8.13.3]{8.13.3}.

\setcounter{subsubsection}{4}\etiquette[8.11.5]{I.sec8.11.5}
\subsubsection{}
现在\pageoriginale 考虑由\sisi{$F$}{$f$} 诱导的函子 $F\mapsto F\circ f$
\hypertarget{sga4.exp.I.equation.8.11.5.1}{}\begin{equation}
f^{\circ\ast}:\check{C}'\longrightarrow \check{C}\tag{8.11.5.1}
\end{equation}
若上述函子将 pro-可表函子映为 pro-可表函子,
则称 $f$ \emph{有一个 pro-伴随}, \ie 存在一个函子
(称为 \emph{$f$ 的 pro-伴随函子})
\hypertarget{sga4.exp.I.equation.8.11.5.2}{}\begin{equation}
g:\Pro(\sisi{C}{C'})\longrightarrow \Pro(\sisi{C'}{C}){\quoi},\tag{8.11.5.2}
\end{equation}
以及一个双函子同构
\hypertarget{sga4.exp.I.equation.8.11.5.3}{}\begin{equation}
\Hom_{\Pro(C)}(g(𝒴'),𝒳)\simeq
\Hom_{\Pro(C')}(𝒴'\sisi{)}{},f(𝒳)){\quoi}.\tag{8.11.5.3}
\end{equation}
当然, 说 $f$ 有 pro-伴随, 就是说 $f^\circ$ 有 ind-伴随,
因此关于 ind-伴随的概念与结果可立即转述为关于 pro-伴随的概念与结果.
这里只指出: $f$ 有 pro-伴随, 当且仅当
$\Pro(f):\Pro(C)\rightarrow\Pro(C')$ 有\emph{\sisi{右}{左}}伴随,
此时上述函子 $g$ 正是 $\Pro(\sisi{C}{f})$ 的一个\sisi{右}{左}伴随;
为此必须 $f$ \emph{左}正合, 而当 $C$ 等价于一个小范畴时这个条件也是充分的.
在这种情形下, $f$ 正合, 当且仅当它同时有 ind-伴随和 pro-伴随.

\begin{enonce*}[remark]{例 8.11.6}\etiquette[8.11.6]{I.exem8.11.6}
考虑函子
$$
f:C\longrightarrow (\Ens){\quoi}.
$$
的情形. 若这个函子有 pro-伴随, 则它 pro-可表;
反之成立, 当且仅当由 pro-可表函子组成的 $\check{C}$ 的全子范畴对小乘积稳定;
特别地, 若 $C$ 等价于小范畴 (\hyperref[I.sec8.10.14]{8.10.14}),
或 $C$ 中小乘积可表 (将\hyperref[I.prop8.9.5]{8.9.5} b) 应用于 $C^\circ$),
则\pageoriginale 属于这一情形.
因此, 除少许保留外, 可以说对一个函子 $f:C\rightarrow C'$,
pro-伴随的存在性是\emph{$f$ 的 pro-可表性概念的自然推广},
后一个概念在 $C'=(\Ens)$ 时有定义.
\end{enonce*}


% END INCLUDED SGA4_English_current_buildable_source_20260906/source/01/drafts/SGA4_I_LINES_5230_5389_English.texfrag


% BEGIN INCLUDED SGA4_English_current_buildable_source_20260906/source/01/drafts/SGA4_I_LINES_5390_5467_English.texfrag
\subsection{严格 ind-对象与 pro-对象. 对可表性判别准则的应用}\etiquette[8.12]{I.sec8.12}

\subsubsection{}\etiquette[8.12.1]{I.sec8.12.1}
设
$$
𝒳=(X_i)_{i\in I}
$$
为 $𝒰$-范畴 $C$ 的一个 ind-对象, $F$ 为它所 ind-表示的预层.
于是立即可见, 说典范态射
$$
X_i\longrightarrow F=\textrm{``}\varinjlim_i\textrm{''\,}X_i
$$
是 $\Ind(C)$ 中的单态射 (或等价地, 是 $\widehat{C}$ 中的单态射,
\ie “逐个变元”看都是函子的单态射),
等价于说对 $I$ 的每个箭头 $i\rightarrow j$, 相应的转移箭头
$$
X_i\longrightarrow X_j
$$
是单态射. 当这些条件满足, 并且 $I$ 是有序范畴时,
称 $𝒳$ 为\emph{严格 ind-对象}.
应\pageoriginale 注意到, 这一条件在 ind-对象的同构下不保持不变;
若一个 ind-对象同构于严格 ind-对象, 则称它\emph{本质严格}.
若一个预层可由严格 ind-对象 ind-表示, 则称它\emph{严格 ind-可表};
因此 $F=\textrm{``}\varinjlim_i\textrm{''\,}X_i$ 严格 ind-可表,
当且仅当 $𝒳=(X_i)_{i\in I}$ 本质严格.

\paragraph{}\etiquette[8.12.1.1]{I.8.12.1.1}
设 $F$ 为 $C$ 上的预层, 考虑 $C_{/F}$ 的如下全子范畴:
它由起点在 $C$ 中且为单态射的箭头 $X\rightarrow F$ 组成.
我们称它为 \emph{$F$ 的可表子函子范畴};
它是 $F$ 的可表子函子所成的有序集对应的范畴,
其序是由 $F$ 的子对象集合的序诱导的序.
于是由定义立即得到:

\begin{enonce*}{命题 8.12.2}\etiquette[8.12.2]{I.prop8.12.2}
$C$ 上的预层 $F$ 严格 ind-可表的充要条件是,
$F$ 的可表子函子的 (有序) 范畴 $I$ 滤过且本质小, 并且有
\hypertarget{sga4.exp.I.equation.8.12.2.1}{}\begin{equation*}
\varinjlim_I X_i \xrightarrow{\sim}F{\quoi}.\tag{8.12.2.1}
\end{equation*}
当对 $C$ 的每个对象 $X$, $C$ 中 $X$ 的子对象集合都是小时
(例如 $C$ 有小的生成子范畴 $\hyperref[I.prop7.4]{(7.4)}$),
这蕴含可表子函子的范畴本身是小的.
\end{enonce*}

\setcounter{subsubsection}{2} \setcounter{paragraph}{1}
\paragraph{}\etiquette[8.12.2.2]{I.sec8.12.2.2}
若\pageoriginale $F$ 严格 ind-可表, 便有一种优先采用的方式用 ind-对象将它 ind-表示,
而且所用的还是严格 ind-对象: 取表示\hyperlink{sga4.exp.I.equation.8.12.2.1}{(8.12.2.1)}.
若严格 ind-对象 $𝒴$ 同构于对某个适当的 $F$,
在\hyperlink{sga4.exp.I.equation.8.12.2.1}{(8.12.2.1)} 中出现的形如 $(X_i)$ 的 ind-对象,
则称它为\emph{饱和的};
因此, 在唯一同构意义下, 与给定的严格 ind-对象同构的饱和严格 ind-对象只有一个,
即在\hyperref[I.prop8.12.2]{8.12.2} 中取
$F=\varinjlim 𝒴$ (在 $\widehat{C}$ 中取极限) 而得到的 ind-对象.

\paragraph{}\etiquette[8.12.2.3]{I.sec8.12.2.3}
设已经知道在集合 $\Ob C_{/F}$ 中能找到一个小的共尾子集;
特别当 $F$ ind-可表时便如此.
则\hyperref[I.prop8.12.2]{8.12.2} 中所考虑的全子范畴 $I$ 也满足同样条件,
所以在判别准则\hyperref[I.prop8.12.2]{8.12.2} 中可以省去 $I$ 本质小的条件.


% END INCLUDED SGA4_English_current_buildable_source_20260906/source/01/drafts/SGA4_I_LINES_5390_5467_English.texfrag


% BEGIN INCLUDED SGA4_English_current_buildable_source_20260906/source/01/drafts/SGA4_I_LINES_5468_5594_English.texfrag
\subsubsection{}\etiquette[8.12.3]{I.sec8.12.3}
设 $F$ 为 $C$ 上的预层, 考虑 $C_{/F}$ 的一个对象
$$
u:X\longrightarrow F
$$
若对 $u$ 的每个形如
\hypertarget{sga4.exp.I.equation.8.12.3.1}{}\begin{equation}
X\xrightarrow{p}X'\xrightarrow{u'}F{\quoi},\tag{8.12.3.1}
\end{equation}
的分解, 其中 $p$ 是严格满态射 (\hyperref[I.sec10.2]{10.2}),
$p$ 都是同构, 则称 $u$ (或对 $(X,u)$) 为\emph{极小的}.
将 $u$ 看作 $F(X)$ 的一个元素\hyperref[I.prop1.4]{(1.4)},
说 $u$ 极小便意味着\pageoriginale 每个满足
$u\in
\Iim(F(p):F(X')\rightarrow F(X))$ 的严格满态射 $p:X\rightarrow X'$
都是同构. 下列引理的 a) 将阐明这一概念:

\begin{enonce*}{引理 8.12.4}\etiquette[8.12.4]{I.lem8.12.4}
设 $F$ 为 $C$ 上的预层.
\begin{enumerate}
\renewcommand{\theenumi}{\alph{enumi}}
\renewcommand{\labelenumi}{\theenumi)}
\item 假定 $F$ 将余核映为核, 并考虑一个态射
$$
u:X\longrightarrow F{\quoi},
$$
其中 $X\in \Ob C$. 当 $C$ 中双箭头的余核可表时,
$u$ 为单态射的充分条件是 $u$ 极小;
若 $C$ 中纤维积可表, 则这个条件也是必要的.

\item 假定 $C$ 中有限 $\varinjlim$ 可表.
要使 $F$ 的可表子函子的子范畴 (\hyperref[I.8.12.1.1]{8.12.1.1})
滤过 (未必小), 且以 $F$ 为归纳极限,
只须 $F$ 左正合, 且每个态射 $u:X\rightarrow F$,
其中 $X\in \Ob C$, 都可分解为
$$
X\rightarrow X'\xrightarrow{u'}F{\quoi},
$$
其中 $u'$ 极小; 若 $C$ 中纤维积可表, 则这个条件也是必要的.

\item 假定 $C$ 有小的生成子范畴 (\hyperref[I.def7.1]{7.1}),
且 $F$ 的可表子函子范畴 $I$ 滤过; 则 $I$ 是小的.
\end{enumerate}
\end{enonce*}

\noindent \emph{证明}.\pageoriginale
\begin{enumerate}
\renewcommand{\theenumi}{\alph{enumi}}
\renewcommand{\labelenumi}{\theenumi)}
\item 假定 $u$ 极小; 来证明 $u$ 是单态射,
\ie 对每对满足 $uv=uv'$ 的箭头 $v,v':Y\rightrightarrows X$, 都有 $v=v'$.
事实上, 若 $X'=\Coker(v,v')$, 则 $u$ 分解为
$X\xrightarrow{p}X'\xrightarrow{u'}F$ (因为 $F$ 将余核映为核),
而由构造 $p$ 为严格满态射, 因此 $p$ 是同构, \ie $v=v'$.
假定 $u$ 是单态射; 来证明它极小.
考虑分解\hyperlink{sga4.exp.I.equation.8.12.3.1}{(8.12.3.1)};
因为 $p$ 为严格满态射, 它是典范双箭头
$v,v':\fprod{X}{X}{X'}\rightrightarrows X$ 的余核,
又因 $(u'p)v=(u'p)v'$ 且 $u'p=u$ 是单态射,
有 $v=v'$, 因此 $p$ 是同构.

\item 充分性: 因为 $F$ 左正合, $C_{/F}$ 滤过.
已知 $F=\varinjlim_{C_{/F}}X$, 由\hyperref[I.prop8.1.3]{8.1.3} c),
只须证明 $C_{/F}$ 中由 $F$ 的子对象组成的全子范畴 $I$ 在 $C_{/F}$ 中共尾;
由 a) 的“充分性”部分, 这正是 $C_{/F}$ 的每个对象都由一个“极小”对象控制这一假设所保证的.
必要性: 因为右正合函子的每个滤过归纳极限仍右正合, 第一个条件显然必要.
第二个条件由 a) 的“必要性”部分得到.

\item 只要知道 $C'_{/F}$ 的全子范畴 $C'_{/X}$,
就知道了 $F$ 的一个可表子函子 $X\hookrightarrow F$,
其中 $C'$ 为 $C$ 的一个固定的小生成子范畴.
(使用 $I$ 滤过这一假设.)
因为 $C'_{/F}$ 是小的, 其全子范畴的集合是小的, 由此即得结论.
\end{enumerate}

\begin{enonce*}{命题 8.12.5}\etiquette[8.12.5]{I.prop8.12.5}
设\pageoriginale $C$ 为 $𝒰$-范畴, 其中有限归纳极限可表,
且它有小的生成子范畴 (\hyperref[I.def7.1]{7.1}).
设 $F$ 为 $C$ 上的预层.
要使 $F$ 严格 ind-可表, 只须它满足下列两个条件;
若 $C$ 中纤维积可表, 则这些条件也是必要的:
\begin{enumerate}
\renewcommand{\theenumi}{\alph{enumi}}
\renewcommand{\labelenumi}{\theenumi)}
\item $F$ 左正合.

\item 每个对 $(X,u)$, 其中 $X\in \Ob C$ 且 $u\in F(X)$,
在 $C_{/F}$ 中都由一个极小对 $\hyperref[I.sec8.12.3]{(8.12.3)}$ 控制,
\ie 存在一个极小对 $(X',u')$ $(X'\in \Ob C, u'\in F(X'))$
以及一个态射 $f:X\rightarrow X'$, 使得 $u=F(f)(u')$.
\end{enumerate}
\end{enonce*}

充分性由\hyperref[I.prop8.12.2]{8.12.2} 及\hyperref[I.lem8.12.4]{8.12.4} b), c) 得到,
必要性由\hyperref[I.prop8.3.1]{8.3.1} 及\hyperref[I.lem8.12.4]{8.12.4} a) 得到.

\begin{enonce*}[remark]{注 8.12.6}\etiquette[8.12.6]{I.rem8.12.6}
若 $F$ 将 $C$ 中的融合和映为纤维积,
则立即可见, 对 b) 中那样一个给定的对 $(X,u)$,
满足 $u\in\Iim(F(X')\rightarrow F(X))$ 的 $X$ 的严格商 $X'$ 的集合向下滤过;
这意味着, 若 $X$ 的严格商的集合是 Artin 的,
则\hyperref[I.prop8.12.5]{8.12.5} 的条件 b) 自动满足.
因此, 若 $C$ 的每个对象 $X$ 都满足上述关于 $X$ 的条件
(此时有时也称 $C^\circ$ 的对象是 \emph{Artin 的}),
则由\hyperref[I.prop8.12.5]{8.12.5} 得知,
$F$ 严格 ind-可表, 当且仅当 $F$ 左正合;
在这种情形下, 只要 $F$ ind-可表, 就严格 ind-可表 (\hyperref[I.prop8.3.1]{8.3.1}).
\end{enonce*}


% END INCLUDED SGA4_English_current_buildable_source_20260906/source/01/drafts/SGA4_I_LINES_5468_5594_English.texfrag


% BEGIN INCLUDED SGA4_English_current_buildable_source_20260906/source/01/drafts/SGA4_I_LINES_5595_5699_English.texfrag
\begin{enonce*}{推论 8.12.7}\etiquette[8.12.7]{I.cor8.12.7}
设\pageoriginale $C$ 为 $𝒰$-范畴, 其中小的投射极限可表,
且它有小的余生成子范畴 (\hyperref[I.sec7.9]{7.9},
\hyperref[I.exem7.13]{7.13}).
则函子 $F:C\rightarrow (\Ens)$ 可表, 当且仅当 $F$ 与小的投射极限可交换.
\end{enonce*}
必要性清楚. 关于充分性, 将\hyperref[I.prop8.12.5]{8.12.5} 应用于对偶范畴 $C^\circ$.
为先证明 $F$ (严格) pro-可表, 还须证明每个对 $(X,u)$,
$X\in \Ob C$ 且 $u\in F(X)$, 都由一个“极小”对控制.
而满足 $u\in\Iim(F(X_i)\rightarrow F(X))$ 的 $X$ 的\emph{严格}子对象族
$(X_i)_{i\in I}$ 是小的 (\hyperref[I.cor7.5]{7.5} 的对偶形式).
因为 $F$ 左正合, 该族余滤过 (\hyperref[I.rem8.12.6]{8.12.6});
又因 $F$ 与小的投射极限可交换, 可见 $X'=\varprojlim_i X_i$ 是这族的最小对象.
(使用 $F$ 左正合, 因而将单态射映为单态射这一事实.)
若 $u'\in F(X')$ 是在 $F(X)$ 中的像为 $u$ 的唯一元素,
则可见 $(X',u')$ 是控制 $(X,u)$ 的一个极小对.
这就证明了 $F$ pro-可表, 然后由下列引理得到结论:

\begin{enonce*}{引理 8.12.8.1}\etiquette[8.12.8.1]{I.lem8.12.8.1}
设 $F:C\rightarrow(\Ens)$ 为函子,
其中 $C$ 是小的投射极限可表的 $𝒰$-范畴.
$F$ 可表的充要条件是它 pro-可表并且与小的投射极限可交换.
\end{enonce*}

必要性\pageoriginale 清楚; 来证明充分性.
说 $F$ 可表, 显然就是说 ${}_F\backslash C$ 有始对象
(事实上, ${}_F\backslash C$ 的始对象恰好是同构 $F\rightarrow X$,
\ie $F$ 的表示数据).
因为 $F$ pro-可表, $({}_F\backslash C)^\circ$ 滤过且等价于一个小范畴,
所以 $X=\varprojlim_{({}_F\backslash C)^\circ} X$ 在 $C$ 中可表;
又因 $F$ 与所考虑的极限可交换, 得知 $X$ 是 ${}_F\backslash C$ 的始对象, \cqfd

\begin{enonce*}{推论 8.12.8}\etiquette[8.12.8]{I.cor8.12.8}
在 $\hyperref[I.cor8.12.7]{8.12.7}$ 关于 $C$ 的假设下,
设 $f:C\rightarrow C'$ 为从 $C$ 到 $𝒰$-范畴 $C'$ 的函子.
$f$ 有左伴随的充要条件是 $f$ 与小的投射极限可交换.
\end{enonce*}

事实上, 将\hyperref[I.cor8.12.7]{8.12.7} 应用于形如
$X\mapsto \Hom(Y',f(X))$ 的复合函子, 便立即归结为该结果.

\begin{enonce*}[remark]{例 8.12.9}\etiquette[8.12.9]{I.exer8.12.9}
如在\hyperref[I.exem7.13]{7.13} 中指出的,
若 $C$ 是拓扑空间 $X\in 𝒰$ (或更一般地, 一个 $𝒰$-景
(II \hyperref[II.def3.0.2]{(3.0.2)}) 上的集合的 $𝒰$-层范畴,
则\hyperref[I.cor8.12.7]{8.12.7} 及\hyperref[I.cor8.12.8]{8.12.8} 中关于 $C$ 的假设满足.
我们给出一个有启发性的例子\footnote[1]{(由 H. BASS 给出).},
说明在\hyperref[I.cor8.12.7]{8.12.7} 中,
$C$ 存在小的余生成子范畴 $D$ 这一假设并非多余.
取 $C$ 为属于 $𝒰$ 的群所成的范畴.
令 $J$ 为 $C$ 中单群的同构类的集合; 对每个 $j\in J$,
在类 $j$ 中选取一个单群 $G_j$;
令 $I$ 为 $J$ 的 $𝒰$-小子集所成的滤过有序集,
对 $i\in I$, 令 $X_i=\amalg_{j\in i}G_j$.
于是各 $X_i$ 组成 $C$ 中的一个\pageoriginale 投射系 $(X_i)_{i\in I}$,
其转移态射为满态射, 而相应的函子
\begin{equation}
F(X)=\varinjlim_i \Hom(X_i,X)\tag{$\ast$}
\end{equation}
取值于 $(𝒰\text{-}\Ens)$
(尽管指标集 $I$ 显然不具有属于 $𝒰$ 的基数);
更精确地, 来证明对 $C$ 的每个小的全子范畴 $C_\circ$,
存在 $i_\circ\in I$, 使得 $F$ 在 $C_\circ$ 上的限制由 $X_{i_\circ}$ 表示;
这将同时证明 $F$ 取值于 $𝒰\text{-}\Ens$, 且与小的投射极限可交换.
要证明这一断言, 只须注意: 对每个 $X\in \Ob C_\circ$,
存在从 $G_j$ 到 $X$ 的非平凡同态的那些 $j\in J$ 所成的集合 $J(X)$ 的基数必小,
因为这样的态射必为单态射 ($G_j$ 是单群);
因此, 若 $i_\circ$ 为对 $X\in\ob C_\circ$ 将各 $J(X)$ 取并而得到的 $J$ 的子集,
则 $i_\circ$ 是小的, \ie $i_\circ\in I$, 且具有所需性质.
另一方面, 显然 $F$ 不可表, 因为 $\card I\not\in 𝒰$.
因此由此及\hyperref[I.cor8.12.7]{8.12.7} 得知:
\emph{$𝒰$ 中的群范畴 $C$ 没有小的全余生成子范畴}.
另一方面, 因为 $C$ 的对象 $\ZZ$ 是一个生成元,
由\hyperref[I.cor7.12]{7.12} 的证明得知,
存在一个双生成群 $G$, 它不能嵌入 $𝒰$ 中的群范畴 $C$ 的任何内射对象.
而且看来 $C$ 很可能除了平凡群外没有内射对象.
\end{enonce*}


% END INCLUDED SGA4_English_current_buildable_source_20260906/source/01/drafts/SGA4_I_LINES_5595_5699_English.texfrag


% BEGIN INCLUDED SGA4_English_current_buildable_source_20260906/source/01/drafts/SGA4_I_LINES_5700_5804_English.texfrag
\subsection{pro-可表函子与可及函子}\etiquette[8.13]{I.sec8.13}\pageoriginale

\subsubsection{}\etiquette[8.13.1]{I.sec8.13.1}
本节使用下一节的一些概念与结果, 特别是\hyperref[I.prop9.11]{9.11} 及
\hyperref[I.prop9.13]{9.13}, 来得到一个 pro-可表性判别准则,
我们将在 IV \hyperref[IV.theo9.5.4]{9.16} 中 (附带地) 使用它.
下文中, $C$ 表示满足\hyperref[I.def9.1]{9.1} b) 的条件 $L$ 的 $𝒰$-范畴;
例如, 若 $C$ 中小的滤过归纳极限可表, 则满足该条件.

\begin{enonce*}{命题 8.13.2}\etiquette[8.13.2]{I.prop8.13.2}
设 $C$ 如上, 且设
$$
f:C\longrightarrow(\Ens)
$$
为一个函子.
\begin{enumerate}
\renewcommand{\theenumi}{\alph{enumi}}
\renewcommand{\labelenumi}{\theenumi)}
\item 假定 $C$ 的每个对象都可及 (\hyperref[I.def9.3]{9.3}).
若 $f$ pro-可表, 则 $f$ 可及 (\hyperref[I.def9.2]{9.2}) 且左正合.

\item 假定 $C$ 中有限投射极限可表, 且 $C$ 有基数滤过
(\hyperref[I.def9.12]{9.12}).
若 $f$ 可及且左正合, 则 $f$ pro-可表.
\end{enumerate}
\end{enonce*}

\noindent \emph{证明.}
\begin{enumerate}
\renewcommand{\theenumi}{\alph{enumi}}
\renewcommand{\labelenumi}{\theenumi)}
\item 关于 $C$ 的假设意味着从 $C$ 到 $(\Ens)$ 的协变可表函子可及.
所以这些函子的每个小的归纳极限也可及 (\hyperref[I.prop9.6]{9.6} (i)),
从而每个 pro-可表函子也可及.

\item 由\hyperref[I.them8.3.3]{8.3.3} (iii),
只须证明 $\Ob({}_F\backslash C)^\circ$ 包含一个小的共尾子范畴.
由假设, 存在基数 $\Pi$, 使得 $f$ 是 $\Pi$-可及的.
设 $(X,u)$,\pageoriginale $u\in
F(X)$, 为 ${}_F\backslash C$ 的一个对象.
使用\hyperref[I.def9.12]{9.12} c) 的记号,
则有 $X=\varinjlim_i X_i$, 其中 $I$ 滤过且相对于 $\Pi$ 足够大,
各 $X_i$ 在 $C'=\Filt^{\Pi}(C)$ 中,
从而 $F(X)\xleftarrow{\sim}\varinjlim_i F(X_i)$.
这表明小子范畴 $({}_F\backslash C')^\circ$
在 $({}_F\backslash C)^\circ$ 中共尾, 证明完成.
\end{enumerate}

\begin{enonce*}{推论 8.13.3}\etiquette[8.13.3]{I.cor8.13.3}
设 $C$ 为满足下列条件的 $𝒰$-范畴:
\begin{enumerate}
\renewcommand{\theenumi}{\alph{enumi}}
\renewcommand{\labelenumi}{\theenumi)}
\item $C$ 中有限投射极限可表.

\item $C$ 中小的滤过归纳极限可表.

\item $C$ 的双箭头范畴上的函子 $\Ker$ 可及
(例如, 它与小的滤过 $\varinjlim$ 可交换).

\item $C$ 的每个严格满态射都是普遍严格满态射 $\hyperref[I.sec10.2]{(10.2)}$.

\item 存在一个通过严格满态射生成的 $C$ 的小子范畴 $\hyperref[I.def7.1]{(7.1)}$.
\end{enumerate}
在这些条件下, 函子 $f:C\rightarrow (\Ens)$ pro-可表,
当且仅当它左正合且可及 (\hyperref[I.def9.2]{9.2}).
\end{enonce*}

事实上, 分别由\hyperref[I.prop9.11]{9.11} 与\hyperref[I.prop9.13]{9.13},
\hyperref[I.prop8.13.2]{8.13.2} a) 与 b) 中关于 $C$ 的条件满足.


\begin{enonce*}{推论 8.13.4}\etiquette[8.13.4]{I.cor8.13.4}
设 $f:C\rightarrow C'$ 为满足\hyperref[I.cor8.13.3]{8.13.3}
的条件 {\rm a)} 至 {\rm e)} 的 $𝒰$-范畴之间的函子.
$f$ 有 pro-伴随的充要条件是 $f$ 左正合且可及.
\end{enonce*}

由于\pageoriginale 可表函子
$h_{Y'}:X'\rightarrow \Hom(Y',X'):C'\rightarrow(\Ens)$ 左正合且可及
(\hyperref[I.prop9.11]{9.11}), 若 $f$ 有同样的性质,
则它与上述函子的复合也有同样性质;
因而由\hyperref[I.cor8.13.3]{8.13.3}, 这些复合 pro-可表,
\ie $f$ 有 pro-伴随.
反之, 假定 $f$ 有 pro-伴随,
令 $C'_1$ 为 $C$ 的一个小的生成子范畴,
$\Pi$ 为一个基数, 使得各 $Y'\in \Ob C'$ 为 $\Pi$-可及;
为证明 $f$ 是 $\Pi$-可及的,
只须证明它与各 $h_{Y'}$, $Y'\in \Ob C'$ 的复合也如此;
由假设, 这些复合 pro-可表, 因而可及 (\hyperref[I.cor8.13.3]{8.13.3}),
所以只要 $\Pi$ 取得充分大, 它们就是 $\Pi$-可及的, \cqfd

% END INCLUDED SGA4_English_current_buildable_source_20260906/source/01/drafts/SGA4_I_LINES_5700_5804_English.texfrag


% BEGIN INCLUDED SGA4_English_current_buildable_source_20260906/source/01/drafts/SGA4_I_LINES_5805_5940_English.texfrag
